\documentclass[12pt, a4paper]{article}

% ── Codificació i idioma ────────────────────────────────────────────────────
\usepackage[utf8]{inputenc}
\usepackage[T1]{fontenc}
\usepackage[spanish]{babel}

% ── Tipografia ──────────────────────────────────────────────────────────────
\usepackage{lmodern}
\usepackage{microtype}

% ── Geometria de la pàgina ──────────────────────────────────────────────────
\usepackage[
  top=2.5cm, bottom=2.5cm,
  left=3cm,  right=3cm
]{geometry}

% ── Capçalera i peu de pàgina ───────────────────────────────────────────────
\usepackage{fancyhdr}
\pagestyle{fancy}
\fancyhf{}
\renewcommand{\headrulewidth}{0.4pt}
\fancyhead[L]{\small Automatització Industrial}
\fancyhead[R]{\small Monitorización de subestaciones eléctricas}
\fancyfoot[C]{\thepage}

% ── Colors ──────────────────────────────────────────────────────────────────
\usepackage[dvipsnames]{xcolor}
\definecolor{upcblue}{RGB}{0, 83, 150}
\definecolor{lightgray}{RGB}{245, 245, 245}
\definecolor{darkgray}{RGB}{80, 80, 80}

% ── Hiperenllaços ───────────────────────────────────────────────────────────
\usepackage[
  colorlinks=true,
  linkcolor=upcblue,
  urlcolor=upcblue
]{hyperref}

% ── Taules ──────────────────────────────────────────────────────────────────
\usepackage{booktabs}
\usepackage{array}
\usepackage{tabularx}

% ── Llistes ─────────────────────────────────────────────────────────────────
\usepackage{enumitem}
\setlist[itemize]{itemsep=3pt, topsep=4pt}

% ── Codi font ───────────────────────────────────────────────────────────────
\usepackage{listings}
\lstset{
  basicstyle=\ttfamily\small,
  backgroundcolor=\color{lightgray},
  frame=single,
  framerule=0.4pt,
  rulecolor=\color{darkgray},
  breaklines=true,
  captionpos=b,
  keywordstyle=\color{upcblue}\bfseries,
  commentstyle=\color{OliveGreen}\itshape,
  stringstyle=\color{BrickRed},
}
\lstdefinelanguage{PHP}{
  morekeywords={new,if,else,echo,function,return,while,foreach,
                isset,empty,die,require,include,class,public,private},
  sensitive=false,
  morecomment=[l]{//},
  morecomment=[s]{/*}{*/},
  morestring=[b]",
  morestring=[b]',
}

% ── Caixes de resultat ──────────────────────────────────────────────────────
\usepackage{tcolorbox}
\tcbuselibrary{skins}
\newtcolorbox{resultbox}[2][]{
  colback=#2!8!white,
  colframe=#2!60!black,
  fonttitle=\bfseries\small,
  title=#1,
  left=6pt, right=6pt, top=4pt, bottom=4pt,
  boxrule=0.5pt,
}

% ── Separació entre paràgrafs ───────────────────────────────────────────────
\setlength{\parskip}{6pt}
\setlength{\parindent}{0pt}

% ════════════════════════════════════════════════════════════════════════════
\begin{document}

% ── Portada ─────────────────────────────────────────────────────────────────
\begin{titlepage}
  \centering
  \vspace*{2cm}
  {\color{upcblue}\rule{\linewidth}{1.5pt}}\\[0.4cm]
  {\Huge\bfseries\color{upcblue} Memòria Descriptiva}\\[0.3cm]
  {\large\color{darkgray} Pràctica: Monitorización de subestaciones eléctricas}\\[0.2cm]
  {\color{upcblue}\rule{\linewidth}{1.5pt}}

  \vspace{1.5cm}

  \begin{tabular}{rl}
    \textbf{Assignatura:}    & Automatització Industrial \\[4pt]
    \textbf{Titulació:}      & Grau en Enginyeria [ESPECIALITAT] \\[4pt]
    \textbf{Escola:}         & [NOM DE L'ESCOLA] -- UPC \\[4pt]
    \textbf{Curs acadèmic:}  & 2025--2026 \\[4pt]
    \textbf{Grup:}           & [NÚM. DE GRUP] \\[4pt]
    \textbf{Data d'entrega:} & [DATA D'ENTREGA] \\
  \end{tabular}

  \vspace{2cm}

  {\large\bfseries Membres de l'equip}\\[0.5cm]
  \begin{tabular}{ll}
    \textbullet\quad [NOM MEMBRE 1] & \texttt{[NIU / correu]} \\[4pt]
    \textbullet\quad [NOM MEMBRE 2] & \texttt{[NIU / correu]} \\[4pt]
    \textbullet\quad [NOM MEMBRE 3] & \texttt{[NIU / correu]} \\
  \end{tabular}

  \vfill
  {\small\color{darkgray} Universitat Politècnica de Catalunya (UPC)}
\end{titlepage}

% ── Índex de continguts ─────────────────────────────────────────────────────
\tableofcontents
\newpage

% ════════════════════════════════════════════════════════════════════════════
\section{Introducció}

L'objectiu d'aquesta pràctica ha estat desenvolupar una aplicació web funcional per a la gestió de sensors en subestacions elèctriques, en el context del monitoratge de dispositius IoT industrials.

L'aplicació s'ha implementat íntegrament amb tecnologies web estàndard: \textbf{PHP natiu} (sense l'ús de cap \textit{framework}), \textbf{HTML5} i \textbf{CSS3 pur}. Per a la capa d'accés a dades s'ha emprat l'extensió \textbf{\texttt{mysqli}} en mode orientat a objectes, instanciada amb:

\begin{lstlisting}[language=PHP, caption={Connexió mysqli orientada a objectes}]
$cn = new mysqli('localhost', 'phpmyadmin', '12345', 'IoTMonitorizacion');
\end{lstlisting}

La base de dades s'anomena \texttt{IoTMonitorizacion} i és accessible a \texttt{localhost} mitjançant \textbf{MariaDB}. L'aplicació permet consultar la llista de sensors existents i inserir-ne de nous, validant les dades d'entrada i gestionant els errors d'integritat de la base de dades.

% ════════════════════════════════════════════════════════════════════════════
\section{Estructura d'Arxius}

L'aplicació web es compon dels cinc fitxers següents:

\subsection{\texttt{estilo.css}}
Fulla d'estils global que defineix l'aparença visual de totes les pàgines. Inclou:
\begin{itemize}
  \item Tipografia de família \textit{sans-serif} (Arial/Helvetica) i contenidor centrat amb ombra.
  \item Format de taules amb \texttt{border-collapse: collapse} i vores en totes les cel·les (\texttt{th}, \texttt{td}).
  \item Estils de formularis, etiquetes i botons diferenciats per acció (enviar / esborrar).
  \item Classes \texttt{.msg-ok} (verd) i \texttt{.msg-error} (vermell) per a la visualització de missatges de resultat.
\end{itemize}

\subsection{\texttt{index-1.html}}
Pàgina principal de l'aplicació. Conté el títol general del projecte i un menú de navegació en forma de llista desordenada (\texttt{<ul>}) amb dos accessos directes:
\begin{itemize}
  \item \textit{Ver lista} $\rightarrow$ \texttt{lista.php}
  \item \textit{Añadir un sensor} $\rightarrow$ \texttt{form-sensor.php}
\end{itemize}

\subsection{\texttt{lista.php}}
Pàgina dinàmica que recupera i mostra tots els sensors registrats. Estableix connexió amb MariaDB, executa un \texttt{SELECT} amb \texttt{INNER JOIN} entre les taules \texttt{Sensor} i \texttt{Magnitud} i presenta els resultats en una taula HTML amb les columnes: \textit{Id, Código, Modelo, Magnitud, Rango}. Allibera la memòria amb \texttt{\$resultado->free()} i tanca la connexió amb \texttt{\$cn->close()}.

\subsection{\texttt{form-sensor.php}}
Formulari d'alta de nous sensors. Es connecta a la base de dades per carregar dinàmicament les magnituds disponibles i poblar un element \texttt{<select>}. Inclou camps de text per a \textit{Código} (amb \texttt{maxlength="6"} i atribut \texttt{required}), \textit{Modelo} i \textit{Rango}, i s'envia per mètode \texttt{GET} cap a \texttt{nuevo-sensor.php}.

\subsection{\texttt{nuevo-sensor.php}}
Script de processament de la inserció. Recull les dades de \texttt{\$\_GET}, valida que els camps obligatoris no estiguin buits i, si la validació és correcta, executa un \texttt{INSERT INTO Sensor} mitjançant un \textbf{prepared statement} amb \texttt{bind\_param} per prevenir injeccions SQL. Gestiona tres estats possibles: dades incompletes, inserció correcta i error de clau duplicada (\texttt{errno~1062} de MariaDB).

% ════════════════════════════════════════════════════════════════════════════
\section{Jocs de Proves}

S'han realitzat tres jocs de proves per verificar el correcte funcionament de la validació i la inserció de dades.

\subsection{Joc de proves 1 --- Enviament de formulari buit}

S'ha accedit a \texttt{form-sensor.php} i s'ha premut el botó \textit{Enviar} sense introduir cap dada. El sistema ha redirigit a \texttt{nuevo-sensor.php}, on s'ha mostrat correctament el missatge d'error:

\begin{resultbox}[Resultat mostrat per pantalla]{BrickRed}
  \texttt{Faltan datos para añadir el sensor}
\end{resultbox}

\textbf{Resultat: CORRECTE.} La validació del camp \textit{Código} (obligatori) ha funcionat com s'esperava.

\subsection{Joc de proves 2 --- Inserció correcta d'un nou sensor}

S'ha completat el formulari amb les dades de la taula~\ref{tab:jp2}:

\begin{table}[h!]
  \centering
  \caption{Dades introduïdes en el joc de proves 2}
  \label{tab:jp2}
  \begin{tabularx}{0.6\textwidth}{>{\bfseries}l X}
    \toprule
    Camp     & Valor introduït   \\
    \midrule
    Magnitud & Temperatura       \\
    Código   & \texttt{APTX14}   \\
    Modelo   & Sensor Test       \\
    Rango    & 0-50 °C           \\
    \bottomrule
  \end{tabularx}
\end{table}

El sistema ha mostrat el missatge:

\begin{resultbox}[Resultat mostrat per pantalla]{OliveGreen}
  \texttt{Operación efectuada correctamente}
\end{resultbox}

Posteriorment, s'ha verificat la inserció consultant la taula \texttt{Sensor} des de \texttt{lista.php}, on el nou registre apareixia correctament.

\textbf{Resultat: CORRECTE.} La inserció a la base de dades s'ha realitzat satisfactòriament.

\subsection{Joc de proves 3 --- Detecció de codi de sensor duplicat}

S'ha tornat a enviar el formulari amb exactament les mateixes dades del joc de proves anterior (Código: \texttt{APTX14}). La taula \texttt{Sensor} té una restricció \texttt{UNIQUE KEY} sobre el camp \texttt{codigo}, de manera que MariaDB ha retornat l'error intern \texttt{errno~1062}. L'aplicació ha capturat aquest error i ha mostrat:

\begin{resultbox}[Resultat mostrat per pantalla]{BrickRed}
  \texttt{Duplicate entry 'APTX14' for key 'codigo'}
\end{resultbox}

\textbf{Resultat: CORRECTE.} El control d'errors d'integritat referencial funciona correctament.

% ════════════════════════════════════════════════════════════════════════════
\section{Repartiment de Tasques}

\begin{table}[h!]
  \centering
  \caption{Distribució de tasques entre els membres del grup}
  \label{tab:tasques}
  \begin{tabularx}{\textwidth}{>{\bfseries}l X}
    \toprule
    Membre            & Tasques realitzades \\
    \midrule
    {[NOM MEMBRE 1]}  & Disseny i configuració de la base de dades local: creació de l'entorn MariaDB, importació de l'esquema SQL (\texttt{IoTMonitorizacion.sql}) i creació de l'usuari de connexió. \\[6pt]
    {[NOM MEMBRE 2]}  & Programació del \textit{backend} PHP: implementació de \texttt{lista.php}, \texttt{form-sensor.php} i \texttt{nuevo-sensor.php}, incloent les connexions \texttt{mysqli} OO, les consultes SQL i la gestió d'errors. \\[6pt]
    {[NOM MEMBRE 1]}  & Disseny de la interfície: creació de \texttt{index-1.html} i \texttt{estilo.css}, maquetació de totes les pàgines i aplicació dels estils visuals. \\[6pt]
    {[NOM MEMBRE 2]}  & Proves d'integració i redacció de la memòria: execució dels tres jocs de proves, verificació dels resultats i elaboració del present document. \\
    \bottomrule
  \end{tabularx}
\end{table}

% ════════════════════════════════════════════════════════════════════════════
\section{Conclusions}

S'ha desenvolupat amb èxit una aplicació web CRUD parcial per a la gestió de sensors industrials. L'ús de \texttt{mysqli} en mode orientat a objectes ha permès establir connexions segures i estructurades amb MariaDB. L'aplicació valida correctament les entrades de l'usuari, gestiona els errors de la base de dades de forma robusta i presenta la informació de manera clara i accessible a través d'una interfície HTML/CSS neta i funcional.

Els tres jocs de proves han confirmat que tots els camins d'execució de l'aplicació (formulari buit, inserció correcta i codi duplicat) funcionen correctament.

\end{document}
